\section{Conclusion}
\label{sec:conclusion}

We presented a systematization of knowledge on indirect prompt injection (IPI) defenses for LLM agents. Using AgentDojo's native end-to-end evaluators---measuring real tool-execution traces for security and real environment state for utility---we evaluated 11 defenses across 4 classes in 1{,}260 trials. Our findings differ materially from proxy-based evaluations: input-encoding defenses (Polymorphic, Mixture-of-Encodings, Spotlighting) achieve the best security-utility tradeoff (ASR 0.14--0.21, Utility 0.67--0.73), reducing attack success by 71--81\% while preserving task completion. Runtime-checking defenses achieve near-zero ASR but sacrifice utility; however, we show that 80\% of this utility loss reflects the base model's capability ceiling rather than defense over-blocking. A prompt-level structured-query approximation without fine-tuning provides little protection (ASR 0.70), confirming that effective structured-query defense requires model training. An ablation of our P2 multi-signal prototype shows the echo-canary provides independent security value (ASR 0.03 alone) but all signal combinations sacrifice utility entirely, refuting the hypothesis that orthogonal multi-signals achieve security at acceptable utility cost. These results demonstrate that end-to-end evaluation with native task-utility and trace-based security metrics is essential for honestly characterizing IPI defense tradeoffs. We release our framework, native evaluation pipeline, and results as artifacts.
